\documentclass[a4paper]{article}
\usepackage{amsfonts}
\usepackage{amsmath}
\usepackage{amssymb}
\usepackage{graphicx}  
\usepackage{cancel}

\newcommand{\degree}{^{\circ}}
\newcommand{\cis}{\operatorname{cis}}
\newcommand{\Bgsin}{\operatorname{Bgsin}}
\newcommand{\Bgcos}{\operatorname{Bgcos}}
\newcommand{\Bgtan}{\operatorname{Bgtan}}
\newcommand{\limit}{\lim\limits}
\newcommand{\summation}{\displaystyle\sum}
\newcommand{\dom}{\operatorname{dom}}
\newcommand{\Int}{\displaystyle\int}

\begin{document}
  
\noindent \large Auteur: Rune Suy \\
\noindent \large Studierichting: Informatica\\
\noindent \large Oefeningenreeks 7A nr. 14.6\\

\medskip

\normalsize

$f(x,y) = 3x + 4y$\\

$\dfrac{\partial f}{\partial x} = 3$\\

$\dfrac{\partial f}{\partial y} = 4$\\

raakvlak $z= f(a) + \dfrac{\partial f}{partial x}(a)(x-a) + \dfrac{\partial f}{\partial y}(a)(y-b)$ met a $= (1,1)$\\

$= 7 + (3x-3) + 4(y-1) = 3x+4y$\\

\begin{thebibliography}{999}
%\bibitem{a} Referentie
\end{thebibliography}
\end{document} 
